\chapter{Probleemanalyse en implementatie van verbeteringen}
\label{ch:verbeteringen}

% \section{Probleemanalyse en implementatie van verbeteringen}
% \label{sec:probleemanalyse-verbeteringen}

Na de nulmeting werd de oorspronkelijke campusapp geanalyseerd op basis van de geautomatiseerde Lighthouse-audit en de manuele toetsenbord- en screenreadercontrole. Uit die analyse bleek dat de belangrijkste knelpunten zich concentreerden rond vier terugkerende thema's: ontbrekende of onduidelijke toegankelijke namen, onvoldoende kleurcontrast, beperkte toetsenbordbediening en ontbrekende labels bij formulier- en selectiecomponenten.

De verbeteringen werden niet willekeurig doorgevoerd, maar stap voor stap geprioriteerd op basis van twee criteria. Ten eerste werd gekeken naar de impact op de gebruikerservaring: problemen die meerdere schermen of veelgebruikte interacties troffen, kregen voorrang. Ten tweede werd rekening gehouden met de technische haalbaarheid, zodat wijzigingen die snel en structureel konden worden doorgevoerd eerst werden aangepakt. Op die manier kon de grootste toegankelijkheidswinst worden behaald binnen de beschikbare implementatietijd.

\subsection{Overzicht van de belangrijkste issues}
\label{subsec:overzicht-belangrijkste-issues}

De nulmeting bracht een aantal terugkerende problemen aan het licht. In de eerste plaats kregen verschillende interactieve elementen geen bruikbare toegankelijke naam mee, waardoor screenreaders ze niet of slechts gedeeltelijk konden aankondigen. Dat geld zowel voor knoppen als voor links en icon-only componenten.

Daarnaast bleek het kleurcontrast op meerdere plaatsen onvoldoende. Vooral teksten op gekleurde achtergronden, statuslabels en bepaalde knoppen zaten onder de aanbevolen contrastdrempel, wat de leesbaarheid bemoeilijkte voor gebruikers met een visuele beperking of een verminderd contrastzicht.

Een derde probleem betrof de toetsenbordnavigatie. Niet alle interactieve elementen waren bereikbaar via het toetsenbord en de focusvolgorde was niet overal logisch of zichtbaar. Hierdoor was de app minder bruikbaar voor gebruikers die geen muis of touchscreen kunnen gebruiken.

Tot slot werden select-elementen en andere invoercomponenten niet consequent gekoppeld aan een zichtbaar of programmatic label. Dat maakte formulieren moeilijker te begrijpen en zorgde voor extra hinder in combinatie met screenreaders.

\subsection{Prioriteiten}
\label{subsec:prioriteiten}

De gevonden problemen werden geordend van hoogste naar laagste prioriteit. Problemen die de kernnavigatie of de belangrijkste interacties blokkeerden, kwamen bovenaan te staan. Daaronder vielen in het bijzonder de knoppen en links zonder toegankelijke naam, omdat die rechtstreeks invloed hebben op de basisbediening van de applicatie.

De tweede prioriteit lag bij de toetsenbordbediening en zichtbare focusstaten. Een applicatie die visueel wel werkt maar via toetsenbord moeilijk bedienbaar is, blijft voor een belangrijke groep gebruikers praktisch ontoegankelijk.

Een derde prioriteit vormden contrastproblemen. Deze issues zijn minder vaak volledig blokkerend, maar hebben wel een duidelijke impact op de leesbaarheid, de begrijpelijkheid en de algemene gebruikservaring.

\begin{table}[h]
\centering
\caption[Prioritering van toegankelijkheidsproblemen]{Overzicht van de belangrijkste toegankelijkheidsproblemen en hun prioriteit.}
\label{tab:prioritering-issues}
\begin{tabular}{p{0.44\textwidth}p{0.18\textwidth}p{0.28\textwidth}}
\hline
\textbf{Issue} & \textbf{Prioriteit} & \textbf{Reden} \\
\hline
Interactieve elementen zonder toegankelijke naam & Hoog & Blokkeren screenreadergebruik en basisnavigatie \\
Toetsenbordfocus en focusvolgorde & Hoog & Essentieel voor niet-muisgebruikers \\
Onvoldoende kleurcontrast & Middel & Vermindert leesbaarheid en visuele toegankelijkheid \\
\hline
\end{tabular}
\end{table}

\subsection{Uitgevoerde verbeteringen}
\label{subsec:uitgevoerde-verbeteringen}

In deze sectie worden de geïdentificeerde toegankelijkheidsproblemen behandeld en de bijhorende verbeteringen toegelicht. 
De verbeteringen worden gegroepeerd per type, zodat duidelijk wordt welke categorie van toegankelijkheidsproblemen werd aangepakt.

Binnen elke categorie wordt telkens één representatief probleem stap voor stap uitgewerkt. 
Daarbij wordt eerst het probleem geanalyseerd, vervolgens de oplossing toegelicht en geïllustreerd met een codefragment (voor en na), eventueel aangevuld met een ondersteunende afbeelding.
Dezelfde aanpak werd nadien consistent toegepast op gelijkaardige elementen binnen de volledige applicatie.

\subsubsection{Toegankelijke namen voor interactieve elementen}

\paragraph{Probleem 1: Ontbrekende toegankelijke naam}

\paragraph{Probleemanalyse}

Tijdens de nulmeting werd vastgesteld dat meerdere interactieve elementen geen toegankelijke naam hadden. 
Dit kwam voornamelijk voor bij icon-gebaseerde knoppen en links/knoppen zonder zichtbare tekst.

Voor schermlezers resulteert dit in aankondigingen zoals \textit{``onbenoemde knop''}, 
waardoor gebruikers niet kunnen afleiden wat de functie van het element is.

\paragraph{Oplossing}

Om dit probleem op te lossen werd gebruikgemaakt van het \texttt{aria-label}-attribuut. 
Dit attribuut voegt een beschrijvende naam toe die door screenreaders wordt voorgelezen.

\begin{listing}[H]
\begin{minted}{csharp}
// Voor
<a href="studentenkaart">
	<i class="fa-regular fa-address-card"></i>
</a>

// Na
<a href="studentenkaart" aria-label="Studentenkaart bekijken">
	<i class="fa-regular fa-address-card"></i>
</a>
\end{minted}
\caption{Toevoegen van een toegankelijke naam via \texttt{aria-label}.}
\label{lst:aria-label}
\end{listing}

\paragraph{Reflectie}

Door het toevoegen van een \texttt{aria-label} krijgt het element een duidelijke betekenis voor screenreaders. 
Hoewel hier slechts één voorbeeld wordt weergegeven, werd deze oplossing systematisch toegepast op alle gelijkaardige interactieve elementen binnen de applicatie.

\paragraph{Probleem 2: Onvoldoende context bij dynamische elementen}

\paragraph{Probleemanalyse}

Naast ontbrekende namen werd vastgesteld dat sommige interactieve elementen wel een naam hadden, maar dat deze onvoldoende context bood of niet specifiek genoeg was. In deze gevallen was een meer beschrijvende, dynamische naamgeving vereist, zoals bij het “volgende les”-element op de homepagina.

\paragraph{Oplossing 1: dynamische aria-labels voor eenvoudige context}

Voor relatief eenvoudige dynamische elementen werd het probleem opgelost door het gebruik van dynamische \texttt{aria-labels}. Hierbij wordt de toegankelijke naam opgebouwd op basis van beschikbare gegevens zoals titel, tijdstip en locatie, zodat de gebruiker een beter beeld krijgt van de inhoud.

\begin{listing}[H]
\begin{minted}{csharp}
// Voor
<a href="/lessenrooster/@Id">
    <div class="class-schedule-block">
        <DateBlock Date="@Start" MonthFormat="ddd" />
        <div class="class-schedule-info">
            <div>
                <p>@Start.ToString("H.mm") - @End.ToString("H.mm")</p>
                <p>@LessonType</p>
            </div>
            <h2>@Name</h2>
            <p>@Room</p>
        </div>
    </div>
</a>

// Na
<a href="/lessenrooster/@Id" aria-label="@($"Volgende les: {Name}. Begint om {Start:H.mm}. Lokaal {Room}.")">
    ...
</a>
\end{minted}
\caption{Gebruik van contextspecifieke aria-labels.}
\label{lst:aria-label-dynamisch}
\end{listing}

\paragraph{Oplossing 2: gecentraliseerde ARIA-logica in dynamische componenten}

Voor complexere dynamische componenten, zoals de nieuws- en eventcarousel op de homepagina, werd gekozen voor een meer gestructureerde en onderhoudbare aanpak. In plaats van de toegankelijke naam inline in de markup op te bouwen, wordt deze centraal gegenereerd via helpermethoden binnen de componentlogica.

Deze aanpak zorgt ervoor dat de ARIA-logica losgekoppeld wordt van de UI-structuur, wat de leesbaarheid en onderhoudbaarheid van de code aanzienlijk verbetert.

\begin{listing}[H]
\begin{minted}{csharp}
<a href="/nieuws-en-events/nieuws/@article.Id"
   class="carousel-item-link"
   aria-label="@GetAriaForNews(article)">
\end{minted}
\caption{Gebruik van een helpermethode voor het genereren van een toegankelijke naam in de nieuws-carousel.}
\label{lst:aria-helper-news}
\end{listing}

De effectieve generatie van de toegankelijke naam gebeurt centraal in de component:

\begin{listing}[H]
\begin{minted}{csharp}
private string GetAriaForNews(NewsDto.Index article)
    => $"Nieuws: {article.Title}";

private string GetAriaForEvent(EventDto.Index article)
{
    var dateText = article.Date.Date
        .ToString("d MMMM", System.Globalization.CultureInfo.CurrentCulture);

    return $"Evenement: {article.Title}. Datum: {dateText}.";
}
\end{minted}
\caption{Centrale generatie van toegankelijke namen voor nieuws- en eventitems.}
\label{lst:aria-helper-methods}
\end{listing}

Door deze aanpak wordt de toegankelijke naam niet langer in de markup opgebouwd, maar centraal beheerd binnen de componentlogica. Dit verhoogt de consistentie van de gegenereerde ARIA-labels en maakt toekomstige aanpassingen eenvoudiger.

\paragraph{Reflectie}

Door het gebruik van dynamische ARIA-labels in combinatie met gecentraliseerde helpermethoden wordt de context van interactieve elementen aanzienlijk verbeterd. Enerzijds wordt snel extra informatie toegevoegd waar nodig, anderzijds wordt in complexere componenten gezorgd voor een schaalbare en onderhoudbare oplossing. Dit resulteert in een betere begrijpbaarheid en voorspelbaarheid voor screenreadergebruikers.

\paragraph{Probleem 3: Onvoldoende communicatie van status en relatie bij uitklapbare elementen}

\paragraph{Probleemanalyse}

Bij uitklapbare of toggle-elementen, zoals een hamburgermenu, werd vastgesteld dat de status (open of gesloten) en de relatie met de bijhorende content niet expliciet gecommuniceerd werden naar screenreaders. 

Hierdoor missen gebruikers essentiële context: het is onduidelijk of een knop een menu opent of sluit en welke inhoud door deze actie beïnvloed wordt. Dit leidt tot verwarring en een minder voorspelbare navigatie-ervaring.

\paragraph{Oplossing}

Om dit probleem te verhelpen werd gebruikgemaakt van de ARIA-attributen \texttt{aria-expanded} en \texttt{aria-controls}. 

Het attribuut \texttt{aria-expanded} geeft de huidige status van het element weer (open of gesloten), terwijl \texttt{aria-controls} een expliciete relatie legt met het element dat door de toggle wordt aangestuurd. 

Bij de implementatie werd rekening gehouden met volgende aandachtspunten:

\begin{itemize}
	\item Het \texttt{id} in \texttt{aria-controls} komt exact overeen met het \texttt{id} van het gecontroleerde element;
	\item De waarde van \texttt{aria-expanded} wordt dynamisch aangepast op basis van de actuele toestand;
	\item De structuur van het DOM blijft consistent zodat screenreaders de relatie correct kunnen interpreteren.
\end{itemize}

\begin{listing}[H]
\begin{minted}{csharp}
@using Microsoft.Extensions.Localization
@using Blazor.WebAssembly.DynamicCulture
@using Rise.Client.Resources.MainComponents.Navigation
@inject IStringLocalizer<Navigation> Loc
@inject IConfiguration Config
@inherits NavBase

<LanguageTrackProvider OnInitializeEvent="provider => provider.RegisterComponent(this)" />

<div>
    <MudIconButton Icon="@CurrentIcon"
                   Color="Color.Inherit"
                   OnClick="ToggleMenu"
                   Class="hamburger-btn" 
                   aria-label="Menu"
                   aria-expanded="@(_isMenuOpen ? "true" : "false")"
                   aria-controls="main-navigation" />
    @if (_isMenuOpen)
    {
        <div class="hamburger-fullscreen" id="main-navigation">
            ...
        </div>
    }
</div>
\end{minted}
\caption{Gebruik van \texttt{aria-expanded} en \texttt{aria-controls} bij een hamburgermenu.}
\label{lst:aria-expanded-controls}
\end{listing}

\paragraph{Reflectie}

Door \texttt{aria-expanded} en \texttt{aria-controls} toe te voegen krijgt een screenreader-gebruiker directe feedback over de toggle-status en de relatie met de inhoud. Dit vermindert verwarring en verbetert de navigatie-ervaring, vooral bij dynamische componenten.

%%%%%%%%%%%%%%%%%%%%%%%%%%%%%%%%%%%%%%%%%%%%%%

\subsubsection{Verbeterde toetsenbordbediening en focus}

Deze verbetering richtte zich op de bedienbaarheid van de applicatie via het toetsenbord. Interactieve elementen werden geanalyseerd op focusbaarheid, volgorde en activering. Elementen die visueel aanwezig waren maar geen focus ontvingen, werden aangepast zodat ze logisch bereikbaar zijn via \texttt{Tab} en \texttt{Shift+Tab}.

Daarnaast werd de zichtbaarheid van de focusstaat verbeterd. Dit is essentieel omdat gebruikers anders moeilijk kunnen inschatten welk element op dat moment actief is. Een duidelijke focusindicator verhoogt de oriëntatie en vermindert het risico op bedieningsfouten.

\paragraph{Probleem: Ontbrekende of onzichtbare focusindicatoren}

\paragraph{Probleemanalyse}

Tijdens de audit werd vastgesteld dat veel interactieve elementen geen duidelijke focusstate hadden. De standaardbrowser-focus (zoals een dunne outline) was op veel plaatsen niet zichtbaar genoeg tegen het achtergrondontwerp van de applicatie. Voor toetsenbordgebruikers is dit zeer problematisch: zij kunnen niet meer volgen welk element op dat moment gefocusseerd is, wat leidt tot desorïentatie en navigatiefouten.

Bovendien waren bepaalde custom-componenten (zoals MudBlazor-knoppen) niet voorzien van enige focusindicator, wat ze praktisch ontoegankelijk maakte voor toetsenbordgebruikers.

\paragraph{Oplossing 1: Centrale focusstijl via \texttt{:focus-visible}}

De oplossing bestaat uit twee delen: ten eerste de CSS-pseudo-klasse \texttt{:focus-visible} gebruiken en toepassen op alle interactieve elementen en ten tweede een duidelijke, consistent zichtbare focusstyle definiëren die goed contrasteert tegen de applicatieachtergrond.

Het \texttt{:focus-visible} pseudo-element zorgt ervoor dat de focusstyle alleen zichtbaar is wanneer het element de focus heeft via het toetsenbord (niet wanneer het via muis geklikt wordt), wat de visuele netheid verbetert.

\begin{listing}[H]
\begin{minted}{css}
a:focus-visible,
button:focus-visible,
input:focus-visible,
select:focus-visible,
textarea:focus-visible,
[mudbla zer-component],
.mud-icon-button:focus-visible,
.mud-button-root:focus-visible,
.mud-menu-button:focus-visible,
[role="link"]:focus-visible,
[tabindex]:focus-visible {
  outline: 3px solid #ff9900;
  outline-offset: 3px;
  border-radius: 6px;
}
\end{minted}
\caption{Centrale CSS-regels voor zichtbare focusindicatoren.}
\label{lst:focus-visible-css}
\end{listing}

De gekozen stijl bestaat uit:

\begin{itemize}
	\item Een \textbf{3px solide oranje outline} (\texttt{\#ff9900}) die goed zichtbaar is tegen lichte en donkere achtergronden;
	\item Een \texttt{outline-offset} van 3px zodat de focuslijn niet direct tegen het element aansluit, wat extra duidelijkheid biedt;
	\item Een \texttt{border-radius} van 6px voor een vriendelijk, consistent uiterlijk.
\end{itemize}

De selectoren dekken alle veelvoorkomende interactieve elementen af:
\begin{itemize}
	\item Native HTML-elementen (\texttt{a}, \texttt{button}, \texttt{input}, \texttt{select}, \texttt{textarea});
	\item MudBlazor-componenten (\texttt{.mud-icon-button}, \texttt{.mud-button-root}, \texttt{.mud-menu-button});
	\item Elementen met ARIA-rollen (\texttt{[role="link"]});
	\item Alle elementen met een \texttt{tabindex} (met opzet focusseerbaar gemaakte elementen).
\end{itemize}

\begin{figure}[H]
\centering

\includegraphics[width=0.38\textwidth]{focus-1-before.PNG}
\caption{Voor de aanpassing: een interactieve element zonder zichtbare focusindicator.}
\label{fig:focus-1-before}
\end{figure}

\begin{figure}[H]
\centering

\includegraphics[width=0.38\textwidth]{focus-1-after.PNG}
\caption{Na de aanpassing: dezelfde interactieve element met een duidelijke oranje focusindicator.}
\label{fig:focus-1-after}
\end{figure}

\paragraph{Reflectie}

Door een expliciete, centraal gedefinieerde focusstyle toe te voegen krijgen toetsenbordgebruikers directe visuele feedback. De keuze voor oranje zorgt voor een duidelijke herkenning zonder de applicatie visueel te overbelasten. Deze ene CSS-regel verbetert de toetsenbordnavigatie aanzienlijk en maakt de applicatie voor alle gebruikers voorspelbaarder.

\paragraph{Oplossing 2: Component-specifieke focusstijl}

\paragraph{Context}

Sommige UI-delen vereisen een afwijkende visuele behandeling om voldoende contrast en herkenbaarheid te garanderen. Een goed voorbeeld is de interactieve campuskaart (SchoonmeersenMap) waar SVG-anchors en kleine klikbare gebieden een strakkere, contrasterende focusindicator nodig hebben.

\paragraph{Implementatie}

Voor dit component werd een component-scoped CSS-regel toegevoegd die een duidelijke rode focusrand en subtiele halo gebruikt. Dit zorgt dat toetsenbordgebruikers de focus ook in complexe SVG-omgevingen duidelijk kunnen volgen.

\begin{listing}[H]
\begin{minted}{css}

.map-svg a:focus-visible,
.map-svg [role="button"]:focus-visible,
.map-svg [tabindex]:focus-visible {
    outline: 3px solid #d32f2f;   /* visible red border */
    outline-offset: 3px;
    box-shadow: 0 0 0 4px rgba(211,47,47,0.12);
}

.map-svg a:focus,
.map-svg [role="button"]:focus,
.map-svg [tabindex]:focus {
    outline: 3px solid #d32f2f !important;
    outline-offset: 3px !important;
}
\end{minted}
\caption{Component-scoped focusregels voor de campuskaart (SchoonmeersenMap).}
\label{lst:map-focus-css}
\end{listing}

\begin{figure}[H]
\centering
\includegraphics[width=0.38\textwidth]{focus-2-before.PNG}
\caption{Voor de aanpassing: interactieve elementen op de campuskaart zonder zichtbare focusindicator.}
\label{fig:focus-2-before}
\end{figure}

\begin{figure}[H]
\centering
\includegraphics[width=0.38\textwidth]{focus-2-after.PNG}
\caption{Na de aanpassing: interactieve elementen op de campuskaart met een duidelijke rode focusindicator.}
\label{fig:focus-2-after}
\end{figure}

\paragraph{Reflectie}

Deze component-specifieke stijl is bedoeld om de leesbaarheid en bediening binnen de kaart te maximaliseren zonder het algemene applicatie-ontwerp te verstoren.

\paragraph{Oplossing 3: Focusbaarheid van niet-interactieve elementen via \texttt{tabindex}}

\paragraph{Probleemanalyse}

Sommige interface-elementen (bijvoorbeeld kaartitems of lesblokken) zijn opgebouwd met semantische maar niet-standaard-interactieve tags zoals \texttt{div} of \texttt{h2}. Wanneer deze elementen klikbaar worden gemaakt via JavaScript of component-events, zijn ze standaard niet bereikbaar met het toetsenbord en dus ontoegankelijk voor toetsenbordgebruikers.

\paragraph{Oplossing}

Een veelgebruikte oplossing is het toevoegen van \texttt{tabindex="0"} om zulke elementen focusbaar te maken. Dit zorgt dat ze opgenomen worden in de tabvolgorde en dat ze visueel een focusstaat krijgen door de eerder beschreven focusregels. Belangrijke aandachtspunten:

\begin{itemize}
    \item Gebruik \texttt{tabindex="0"} (niet positieve waarden) zodat de natuurlijke tabvolgorde behouden blijft;
    \item Geef duidelijke ARIA- of tekstuele labels (\texttt{aria-label} of zichtbare tekst) zodat screenreaders context krijgen;
    \item Zorg dat toetsenbordactivatie werkt: behandel \texttt{Enter} en \texttt{Space} als activering (bij voorkeur via een aparte \texttt{@onkeydown}-handler of door een semantisch element zoals \texttt{button} of \texttt{a} te gebruiken);
    \item Overweeg \texttt{role="button"} wanneer een element semantisch geen knop is maar wel knopgedrag heeft.
\end{itemize}
Onderstaand voorbeeld illustreert de aanpak toegepast op een lesblok in het lessenrooster. Het element is focusbaar gemaakt met \texttt{tabindex="0"} en krijgt een dynamische \texttt{aria-label} zodat screenreaders relevante context voorlezen.

\begin{listing}[H]
\begin{minted}{csharp}
//voor
<div aria-label="@($"les: {Course.Name}, Begint om {Course.Start:HH.mm}.")" 
     class="class-schedule-block"
     @onclick="() => OnCourseSelected.InvokeAsync(Course!.Id)"
     @onkeydown="@(e => { if (e.Key == "Enter" || e.Key == " ") OnCourseSelected.InvokeAsync(Course!.Id); })">
    @if (Course is not null)
    {
        <TimeBlock StartTime="@Course.Start" EndTime="@Course.End"/>
        <div class="class-schedule-info">
            <h2 class="course-name">@Course.Name</h2>
            ...
        </div>
    }
</div>

//Na
<div tabindex="0" aria-label="@($"les: {Course.Name}, Begint om {Course.Start:HH.mm}.")" 
     class="class-schedule-block"
     @onclick="() => OnCourseSelected.InvokeAsync(Course!.Id)"
     @onkeydown="@(e => { if (e.Key == "Enter" || e.Key == " ") OnCourseSelected.InvokeAsync(Course!.Id); })">
    @if (Course is not null)
    {
        <TimeBlock StartTime="@Course.Start" EndTime="@Course.End"/>
        <div class="class-schedule-info">
            <h2 class="course-name">@Course.Name</h2>
            ...
        </div>
    }
</div>
\end{minted}
\caption{Gebruik van \texttt{tabindex="0"} en keyboard-activatie op een lesblok.}
\label{lst:tabindex-course}
\end{listing}

\begin{figure}[H]
\centering
\includegraphics[width=0.38\textwidth]{focus-3-before.PNG}
\caption{Voor de aanpassing: niet-interactieve elementen zonder zichtbare focusindicator.}
\label{fig:focus-3-before}
\end{figure}

\begin{figure}[H]
\centering
\includegraphics[width=0.38\textwidth]{focus-3-after.PNG}
\caption{Na de aanpassing: niet-interactieve elementen en zichtbare focusindicator.}
\label{fig:focus-3-after}
\end{figure}

\paragraph{Reflectie}

Het gebruik van \texttt{tabindex} maakt dynamische of niet-semantische elementen direct toegankelijk voor toetsenbordgebruikers. Waar mogelijk verdient het de voorkeur om semantische interactieve elementen (\texttt{button}, \texttt;a\texttt{}) te gebruiken; \texttt{tabindex} is echter een pragmatische en veilige oplossing wanneer herstructurering niet haalbaar is. Zorg altijd voor bijbehorende keyboard-handlers en ARIA-annotaties om volledige toegankelijkheid te garanderen.

\paragraph{Integrale toepassing van focusverbeteringen}

De beschreven oplossingen werden niet beperkt tot individuele componenten, maar systematisch toegepast over de volledige applicatie. Door het combineren van een centrale focusstijl met component-specifieke uitbreidingen en het verbeteren van focusbaarheid via \texttt{tabindex}, werd een consistente en voorspelbare toetsenbordnavigatie gerealiseerd.

Deze geïntegreerde aanpak zorgt ervoor dat gebruikers zich vlot door de interface kunnen bewegen, ongeacht het type component of complexiteit van de interactie. Hierdoor wordt de toegankelijkheid niet enkel lokaal verbeterd, maar structureel verankerd binnen de applicatie.

%%%%%%%%%%%%%%%%%%%%%%%%%%%%%%%%%%%%%%%%%%%%%%

\subsubsection{Contrast en visuele leesbaarheid}

Een derde cluster van verbeteringen betreft het kleurcontrast. Teksten, iconen en belangrijke interface-elementen worden herbekeken zodat ze beter leesbaar zijn tegen hun achtergrond. Waar nodig worden kleuren bijgestuurd of wordt de visuele hiërarchie aangepast.

Deze wijzigingen dragen bij aan de waarneembaarheid van de inhoud en zijn vooral relevant voor gebruikers met een visuele beperking (kleurenblindheid of tijdelijke visuele beperkingen), bijvoorbeeld bij gebruik in felle omgevingen.

\paragraph{Probleem: Onvoldoende kleurcontrast}

\paragraph{Probleemanalyse}

Tijdens de Lighthouse-audit en visuele beoordeling werd opgemerkt dat bepaalde tekstkleuken onvoldoende contrast hadden tegen hun achtergrond. Dit was vooral zichtbaar bij:

\begin{itemize}
	\item Statuslabels (bijvoorbeeld deadlines) met donkergrijze of bruine tinten (zoals \texttt{\#DBDBC1});
	\item Nieuwskaart-datums met niet-contrastrijke kleuren (zoals \texttt{\#bd5b5b});
	\item Andere metadata-elementen met grijze tinten tegen lichte achtergronden.
\end{itemize}

Deze kleuren zorgden ervoor dat teksten moeilijk leesbaar waren, vooral voor gebruikers met een verminderd zicht of bij gebruik in omgevingen met fel licht.

\paragraph{Oplossing 1: Kleuren van elementaire componenten aanpassen}

De eerste stap was het systematisch herzien en aanpassen van de CSS-kleurwaarden in componenten. Voor elementen met onvoldoende contrast werden de kleuren vervangen door varianten met beter contrast. Bijvoorbeeld:

\begin{itemize}
	\item Deadline-labels: grijze tint vervangen door een contrast-rijker alternatief;
	\item Nieuwskaart-datums: rozeachtige tint (\texttt{\#bd5b5b}) geoptimaliseerd;
	\item Cursustype-labels: grijze tinten vervangen.
\end{itemize}

Deze wijzigingen werden doorgevoerd in de component-scoped CSS-bestanden:

\begin{listing}[H]
\begin{minted}{css}
/* VOORDIEN: onvoldoende contrast */
.deadline-info p {
  color: #DBDBC1;
  margin-bottom: 0;
}

.news-card-date {
  font-size: 0.85rem;
  color: #bd5b5b;
}

.class-schedule-info p:first-of-type {
  color: #DBDBC1; 
}

/* NADIEN: verbeterd contrast */
.deadline-info p {
  color: #000000;
  margin-bottom: 0;
}

.news-card-date {
  font-size: 0.85rem;
  color: #000000;
}

.class-schedule-info p:first-of-type {
  color: #000000;
}
\end{minted}
\caption{Voorbeelden van contrastverbeteringen in component-CSS.}
\label{lst:contrast-colors}
\end{listing}

\begin{figure}[H]
\centering
\includegraphics[width=0.75\textwidth]{homepage-contrast-before.png}
\caption{Voor de aanpassing: volgende les component met lage contrastwaarde.}
\label{fig:next-lesson-contrast-before}
\end{figure}

\begin{figure}[H]
\centering
\includegraphics[width=0.75\textwidth]{homepage-contrast-after.png}
\caption{Na de aanpassing: volgende les component met verbeterde contrastwaarde.}
\label{fig:next-lesson-contrast-after}
\end{figure}

\paragraph{Oplossing 2: MudBlazor-componentkleuren aanpassen}

Naast eigen componenten werden ook standaardcomponenten van MudBlazor aangepast. MudBlazor stelt standaard kleuren in die niet altijd optimaal contrast hebben. Een specifiek voorbeeld is het \texttt{.mud-alert-message} element, waarvan de standaard tekstkleur werd aangepast:

Een concreet voorbeeld is de aanpassing van de tekstkleur binnen alertcomponenten:

\begin{listing}[H]
\begin{minted}{css}
/* Alert-bericht kleur aanpassen voor beter contrast */
.alert-section ::deep .mud-alert-message {
  color: #1c4bba;  /* MUD-blauw, contrastrijker dan standaard */
}
\end{minted}
\caption{Aanpassing van MudBlazor alert-kleur voor verbeterd contrast.}
\label{lst:mud-alert-contrast}
\end{listing}

Door deze overschrijvingen wordt de standaard styling van MudBlazor aangepast naar een versie die beter aansluit bij de contrastvereisten van de applicatie.

\begin{figure}[H]
\centering
\includegraphics[width=0.75\textwidth]{homepage-alert-before.png}
\caption{Voor de aanpassing: alertweergave met lage contrastwaarde.}
\label{fig:alert-contrast-before}
\end{figure}

\begin{figure}[H]
\centering
\includegraphics[width=0.75\textwidth]{homepage-alert-after.png}
\caption{Na de aanpassing: alertweergave met verbeterde contrastwaarde.}
\label{fig:alert-contrast-after}
\end{figure}

\paragraph{Reflectie}

De combinatie van component-level kleurverbeteringen en MudBlazor-overschrijvingen zorgt voor een consistent contrastpatroon doorheen de interface. Dit verbetert niet alleen de leesbaarheid voor gebruikers met visuele beperkingen, maar verhoogt ook de algehele leesbaarheid en scherpheid van de gebruikersinterface.

%%%%%%%%%%%%%%%%%%%%%%%%%%%%%%%%%%%%%%%%%%%%%%%%%%%

\subsection{Effect op toegankelijkheid en kwaliteit}
\label{subsec:effect-toegankelijkheid-kwaliteit}

De uitgevoerde verbeteringen hebben een directe impact op de dagelijkse bruikbaarheid van de campusapp. Gebruikers die vertrouwen op toetsenbordnavigatie of een screenreader kunnen nu een groter deel van de applicatie zelfstandig bedienen. Tegelijk verhoogt een consistente focus, duidelijke naamgeving en beter contrast ook de gebruiksvriendelijkheid voor gebruikers zonder specifieke beperking.

De implementatie toont daarnaast aan dat toegankelijkheidsmaatregelen niet noodzakelijk een negatieve invloed hebben op de technische kwaliteit van de applicatie. Integendeel: betere semantiek, duidelijkere componentstructuur en minder ambiguïteit in de interface dragen vaak bij aan een robuustere codebasis.

\subsection{Tussentijdse reflectie}
\label{subsec:tussentijdse-reflectie}

Tijdens de implementatie bleek dat sommige toegankelijkheidsproblemen relatief eenvoudig op componentniveau konden worden opgelost, terwijl andere issues een bredere herwerking van interactiepatronen vereisten. Vooral de structurele problemen rond focusbeheer en interactieve labels vroegen meer aandacht dan de visuele aanpassingen.

De gekozen aanpak bevestigt het belang van prioritering. Door eerst de meest impactvolle drempels aan te pakken, kon de applicatie stap voor stap toegankelijker worden gemaakt zonder onnodig veel tijd te investeren in minder urgente optimalisaties. Deze tussenfase vormt daarmee de brug tussen de initiële analyse en de hertest in het volgende hoofdstuk.
